\begin{abstract}

The increasing demand for efficient and transparent electricity management has highlighted the limitations of conventional postpaid and centralized electricity systems. These systems often provide limited real-time information about energy consumption and depend heavily on centralized mechanisms for billing and account management. This project proposes a software-based \textbf{Decentralized Prepaid Smart Power Grid} that combines blockchain-based prepaid electricity management, machine learning, simulated energy telemetry, and web-based monitoring within a hybrid system architecture.

The proposed system uses a blockchain smart contract to manage prepaid electricity tokens, user balances, and authorization states. Simulated electricity consumption telemetry is processed through a FastAPI backend and stored in a SQLite database. Machine learning techniques are incorporated to provide intelligent analysis of consumption data, with the Isolation Forest algorithm used for anomaly detection and Linear Regression used for short-term consumption forecasting. The system also estimates the remaining usage duration based on the available prepaid balance and estimated consumption rate. A React-based web interface provides users with information regarding electricity consumption, prepaid balance, forecasted usage, authorization status, and detected anomalies.

The proposed architecture is hybrid in nature. Blockchain provides the decentralized component for prepaid balance and authorization management, while the backend, database, machine-learning module, and frontend provide centralized application services. Physical electricity meters, sensors, relays, and power-distribution hardware are outside the scope of the current project; therefore, electricity consumption is represented through simulated telemetry.

The project is expected to demonstrate the feasibility of integrating blockchain-based prepaid management with machine-learning-based consumption analysis and web-based monitoring. The resulting prototype can serve as a foundation for future integration with physical smart-metering infrastructure, real electrical load control, and deployment on a production blockchain network.

\end{abstract}